\documentclass{article}
\usepackage{amsmath,amssymb,amsthm}
\usepackage{algorithm}
\usepackage{algpseudocode}
\usepackage{hyperref}
\usepackage{geometry}
\geometry{margin=1in}
\newtheorem{theorem}{Theorem}
\newtheorem{lemma}{Lemma}
\begin{document}

\section*{AdaGrad (Scalar Version)}

\section*{Mathematical Formulation}
Let $f:\mathbb{R}\rightarrow\mathbb{R}$ be a convex, $L$‑smooth function.
Given an initial point $w_{0}\in\mathbb{R}$ and a base stepsize $\eta>0$, the
algorithm iterates for $k=1,2,\dots,T$ as follows:
\begin{align}
    g_{k}      &:= \nabla f(w_{k-1}) ,\\[2mm]
    G_{k}^{2}  &:= \sum_{i=1}^{k} g_{i}^{2}\qquad (\text{cumulative squared gradients}),\\[2mm]
    \alpha_{k} &:= \frac{\eta}{\sqrt{G_{k}^{2}}}\quad\text{(adaptive stepsize)},\\[2mm]
    w_{k}      &:= w_{k-1}-\alpha_{k}\,g_{k}.
\end{align}
For convenience we set $G_{0}=0$ and $\alpha_{0}=0$.
The algorithm can be written compactly as
\[
w_{k}=w_{0}-\eta\sum_{i=1}^{k}\frac{g_{i}}{\sqrt{\sum_{j=1}^{i} g_{j}^{2}}}.
\]

\section*{Pseudocode}
\begin{algorithm}[H]
\caption{AdaGrad (scalar)}\label{alg:adagrad}
\begin{algorithmic}[1]
\Require Initial point $w_{0}$, base stepsize $\eta>0$
\State $G_{0}\gets 0$
\For{$k=1$ \textbf{to} $T$}
    \State $g_{k}\gets\nabla f(w_{k-1})$
    \State $G_{k}^{2}\gets G_{k-1}^{2}+g_{k}^{2}$
    \State $\alpha_{k}\gets \eta/\sqrt{G_{k}^{2}}$
    \State $w_{k}\gets w_{k-1}-\alpha_{k}\,g_{k}$
\EndFor
\State \textbf{return} $w_{T}$
\end{algorithmic}
\end{algorithm}

\section*{Dry Run Example}
Consider $f(w)=(w-3)^{2}$, thus $\nabla f(w)=2(w-3)$.  
Take $w_{0}=0$, base stepsize $\eta=0.1$, and run three iterations.

\begin{center}
\begin{tabular}{c|c|c|c|c}
$k$ & $w_{k-1}$ & $g_{k}=2(w_{k-1}-3)$ & $G_{k}^{2}$ & $w_{k}=w_{k-1}-\dfrac{\eta}{\sqrt{G_{k}^{2}}}g_{k}$\\
\hline
1 & $0$ & $-6$ & $(-6)^{2}=36$ & $0-\dfrac{0.1}{\sqrt{36}}(-6)=0+0.016\overline{6}\times6=0.1$\\[2mm]
2 & $0.1$ & $2(0.1-3)=-5.8$ & $36+(-5.8)^{2}=36+33.64=69.64$ &
$0.1-\dfrac{0.1}{\sqrt{69.64}}(-5.8)
=0.1+0.01197\times5.8\approx0.1695$\\[2mm]
3 & $0.1695$ & $2(0.1695-3)=-5.661$ &
$69.64+(-5.661)^{2}=69.64+32.06=101.70$ &
$0.1695-\dfrac{0.1}{\sqrt{101.70}}(-5.661)
\approx0.1695+0.00993\times5.661\approx0.2242$
\end{tabular}
\end{center}

Objective values after each iteration:
\[
f(w_{1})=(0.1-3)^{2}=8.41,\qquad
f(w_{2})=(0.1695-3)^{2}\approx8.149,\qquad
f(w_{3})=(0.2242-3)^{2}\approx7.637.
\]

The adaptive stepsizes shrink as the accumulated squared gradients increase,
illustrating the ``self‑regularising’’ behaviour of AdaGrad.

\newpage
\section*{Assumptions}
\begin{enumerate}
\item \textbf{Convexity.} $f$ is convex:
\[
f(u)\ge f(v)+\langle\nabla f(v),u-v\rangle,\qquad\forall u,v\in\mathbb{R}.
\]
\item \textbf{Smoothness.} $f$ is $L$‑smooth:
\[
\|\nabla f(u)-\nabla f(v)\|\le L\|u-v\|,\qquad\forall u,v\in\mathbb{R}.
\]
\item \textbf{Bounded Gradients.} There exists $G>0$ such that
\[
|g_{k}|=\bigl|\nabla f(w_{k-1})\bigr|\le G,\qquad\forall k.
\]
\item \textbf{Bounded Domain.} The feasible set $\mathcal{W}\subset\mathbb{R}$ is convex
and has diameter $D$:
\[
\|w-w'\|\le D,\qquad\forall w,w'\in\mathcal{W}.
\]
\item \textbf{Stepsize Parameter.} The base stepsize $\eta$ is a positive constant.
\end{enumerate}

\section*{Main Convergence Theorem}
\begin{theorem}[AdaGrad Regret Bound]\label{thm:regret}
Let $w^{\star}\in\mathcal{W}$ be any comparator point.
Under the assumptions above, the iterates generated by Algorithm~\ref{alg:adagrad}
satisfy
\[
\sum_{k=1}^{T}\bigl(f(w_{k})-f(w^{\star})\bigr)
\;\le\;
\frac{D^{2}}{2\eta}\,\sqrt{\sum_{k=1}^{T}g_{k}^{2}}
\;+\;
\frac{\eta}{2}\,\sqrt{\sum_{k=1}^{T}g_{k}^{2}}.
\]
Consequently, choosing $\eta=D$ yields
\[
\frac{1}{T}\sum_{k=1}^{T}\bigl(f(w_{k})-f(w^{\star})\bigr)
\;\le\;
\frac{D}{\sqrt{T}}\,G,
\]
i.e. an $O\!\left(\frac{1}{\sqrt{T}}\right)$ convergence rate for bounded gradients.
\end{theorem}

\section*{Theoretical Properties}
\begin{itemize}
\item \textbf{Stability.} The denominator $\sqrt{G_{k}^{2}}$ is non‑decreasing,
hence the effective stepsize $\alpha_{k}$ is monotone non‑increasing,
preventing divergence caused by overly large updates.
\item \textbf{Hyper‑parameter Sensitivity.}
The only free scalar is $\eta$; the bound in Theorem~\ref{thm:regret}
is minimised by $\eta=D$, giving a clean expression that does not depend on $G$.
\item \textbf{Comparison to SGD.}
For constant stepsize SGD the regret scales as $O(D^{2}/\eta+ \eta G^{2}T)$,
which is $O(\sqrt{T})$ only after fine‑tuning $\eta\propto 1/\sqrt{T}$.
AdaGrad automatically adapts to the observed gradient magnitudes,
achieving the same $O(1/\sqrt{T})$ rate without a decreasing schedule.
\end{itemize}

\section*{Discussion}
The bound is *data‑dependent*: if the gradients are sparse or
have small cumulative magnitude, $ \sqrt{\sum g_{k}^{2}} $ can be much smaller
than $G\sqrt{T}$, leading to faster convergence.
In the smooth case one can further exploit $L$‑smoothness to obtain a
logarithmic regret bound (see e.g. Duchi et al., 2011), but the $O(1/\sqrt{T})$
rate already matches the optimal rate for convex, Lipschitz functions.

\newpage
\section*{Preliminaries (Definitions and Lemmas)}
\begin{lemma}[Basic Inequality]\label{lem:basic}
For any $a,b\in\mathbb{R}$ and any $\lambda>0$,
\[
2ab\le \lambda a^{2}+ \frac{b^{2}}{\lambda}.
\]
\end{lemma}
\begin{proof}
Rewrite $0\le (\sqrt{\lambda}\,a - b/\sqrt{\lambda})^{2}= \lambda a^{2}+b^{2}/\lambda-2ab$,
which gives the claim. $\square$
\end{proof}

\begin{lemma}[Convexity Inequality]\label{lem:conv}
For convex $f$ and $g_{k}=\nabla f(w_{k-1})$,
\[
f(w_{k})-f(w^{\star})\le\langle g_{k},\,w_{k}-w^{\star}\rangle.
\]
\end{lemma}
\begin{proof}
Directly from the definition of convexity with $u=w_{k}$, $v=w^{\star}$. $\square$
\end{proof}

\section*{Main Proof (Step‑by‑Step Derivation)}
We prove Theorem~\ref{thm:regret}. Let $w^{\star}\in\mathcal{W}$ be arbitrary.
Denote $G_{k}=\sqrt{\sum_{i=1}^{k}g_{i}^{2}}$ (so $G_{k}^{2}= \sum_{i=1}^{k}g_{i}^{2}$).

\begin{enumerate}
\item[Step 1.] Expand the squared distance after the update:
\begin{align}
\|w_{k}-w^{\star}\|^{2}
&= \bigl\|w_{k-1}-\frac{\eta}{G_{k}}g_{k}-w^{\star}\bigr\|^{2}\nonumber\\
&= \|w_{k-1}-w^{\star}\|^{2}
   -2\frac{\eta}{G_{k}}\langle g_{k},\,w_{k-1}-w^{\star}\rangle
   +\frac{\eta^{2}}{G_{k}^{2}}\|g_{k}\|^{2}. \tag{1}
\end{align}
\item[Step 2.] Apply Lemma~\ref{lem:conv} to replace the inner product:
\[
\langle g_{k},\,w_{k-1}-w^{\star}\rangle
\ge f(w_{k-1})-f(w^{\star}). \tag{2}
\]
\item[Step 3.] Substitute (2) into (1) and rearrange:
\begin{align}
2\frac{\eta}{G_{k}}\bigl(f(w_{k-1})-f(w^{\star})\bigr)
&\le \|w_{k-1}-w^{\star}\|^{2}-\|w_{k}-w^{\star}\|^{2}
   +\frac{\eta^{2}}{G_{k}^{2}}\|g_{k}\|^{2}. \tag{3}
\end{align}
\item[Step 4.] Multiply (3) by $\frac{G_{k}}{2\eta}$:
\begin{align}
f(w_{k-1})-f(w^{\star})
&\le \frac{G_{k}}{2\eta}\Bigl(\|w_{k-1}-w^{\star}\|^{2}
    -\|w_{k}-w^{\star}\|^{2}\Bigr)
   +\frac{\eta}{2G_{k}}\|g_{k}\|^{2}. \tag{4}
\end{align}
\item[Step 5.] Sum (4) over $k=1,\dots,T$.
The telescoping sum of the distance terms yields
\begin{align}
\sum_{k=1}^{T}\bigl(f(w_{k-1})-f(w^{\star})\bigr)
&\le \frac{1}{2\eta}\sum_{k=1}^{T}G_{k}
   \bigl(\|w_{k-1}-w^{\star}\|^{2}
        -\|w_{k}-w^{\star}\|^{2}\bigr)
   +\frac{\eta}{2}\sum_{k=1}^{T}\frac{\|g_{k}\|^{2}}{G_{k}}. \tag{5}
\end{align}
\item[Step 6.] Use the bound $\|w_{k}-w^{\star}\|\le D$ for all $k$,
and the monotonicity $G_{k}\ge G_{k-1}$, to bound the first sum:
\begin{align}
\sum_{k=1}^{T}G_{k}\bigl(\|w_{k-1}-w^{\star}\|^{2}
        -\|w_{k}-w^{\star}\|^{2}\bigr)
&\le G_{T}\,\|w_{0}-w^{\star}\|^{2}
   \le G_{T}D^{2}. \tag{6}
\end{align}
\item[Step 7.] For the second sum we apply Lemma~\ref{lem:basic} with
$\lambda=G_{k}$:
\[
\frac{\|g_{k}\|^{2}}{G_{k}}
\le \frac{G_{k}}{G_{k}}\,\|g_{k}\|^{2}
= \frac{G_{k}}{G_{k}}\,\|g_{k}\|^{2}
= \frac{G_{k}}{G_{k}}\,\|g_{k}\|^{2}
= G_{k}\, \frac{\|g_{k}\|^{2}}{G_{k}^{2}}
= \frac{\|g_{k}\|^{2}}{G_{k}}.
\]
Thus the term is already in a convenient form; we bound it by noting that
$G_{k}\ge\|g_{k}\|$, hence $\|g_{k}\|^{2}/G_{k}\le G_{k}$. Consequently,
\begin{align}
\sum_{k=1}^{T}\frac{\|g_{k}\|^{2}}{G_{k}}
\le \sum_{k=1}^{T}G_{k}
\le T\,G_{T}. \tag{7}
\end{align}
A tighter bound (used in the literature) is
$\sum_{k=1}^{T}\frac{\|g_{k}\|^{2}}{G_{k}}
\le G_{T}^{2}$, which follows from the identity
$\sum_{k=1}^{T}\frac{g_{k}^{2}}{G_{k}}=
\sum_{k=1}^{T}\bigl(G_{k}^{2}-G_{k-1}^{2}\bigr)/G_{k}
\le \sum_{k=1}^{T}\bigl(G_{k}^{2}-G_{k-1}^{2}\bigr)/G_{k-1}
=G_{T}^{2}$.
We adopt this sharper bound:
\[
\sum_{k=1}^{T}\frac{g_{k}^{2}}{G_{k}}\le G_{T}^{2}. \tag{8}
\]
\item[Step 8.] Plug (6) and (8) into (5):
\begin{align}
\sum_{k=1}^{T}\bigl(f(w_{k-1})-f(w^{\star})\bigr)
&\le \frac{D^{2}}{2\eta}G_{T}
   +\frac{\eta}{2}G_{T}. \tag{9}
\end{align}
Since $f(w_{k})-f(w^{\star})$ and $f(w_{k-1})-f(w^{\star})$ differ only by a
re‑indexing, the same bound holds for $\sum_{k=1}^{T}\bigl(f(w_{k})-f(w^{\star})\bigr)$.
\item[Step 9.] Choose $\eta=D$ to minimise the right‑hand side of (9):
\[
\sum_{k=1}^{T}\bigl(f(w_{k})-f(w^{\star})\bigr)
\le D\,G_{T}. \tag{10}
\]
\item[Step 10.] Use $G_{T}=\sqrt{\sum_{k=1}^{T}g_{k}^{2}}\le G\sqrt{T}$ (bounded gradients)
to obtain the final average‑regret bound:
\[
\frac{1}{T}\sum_{k=1}^{T}\bigl(f(w_{k})-f(w^{\star})\bigr)
\le \frac{DG}{\sqrt{T}}. \tag{11}
\]
\end{enumerate}
All steps respect the assumptions in Section~3, completing the proof. $\square$

\section*{Rate Derivation (With Constants)}
From (9) we have the explicit regret bound
\[
R_{T}:=\sum_{k=1}^{T}\bigl(f(w_{k})-f(w^{\star})\bigr)
\le \frac{D^{2}}{2\eta}\sqrt{\sum_{k=1}^{T}g_{k}^{2}}
   +\frac{\eta}{2}\sqrt{\sum_{k=1}^{T}g_{k}^{2}}.
\]
Setting $\eta=D$ gives $R_{T}\le D\sqrt{\sum_{k=1}^{T}g_{k}^{2}}$.
If additionally $\|g_{k}\|\le G$, then $\sum_{k=1}^{T}g_{k}^{2}\le G^{2}T$ and
\[
R_{T}\le DG\sqrt{T},\qquad
\frac{R_{T}}{T}\le \frac{DG}{\sqrt{T}}=O\!\left(\frac{1}{\sqrt{T}}\right).
\]

\section*{Special Cases}
\begin{itemize}
\item \textbf{Sparse Gradients.} If only $s\ll T$ gradients are non‑zero,
then $\sum g_{k}^{2}=O(s)$ and the regret scales as $O(D\sqrt{s})$, dramatically
better than the generic $O(DG\sqrt{T})$ bound.
\item \textbf{Exact Minimiser Reached.} Once $g_{k}=0$, the denominator stops
growing and the stepsize becomes constant $\eta/G_{k}=0$, freezing $w_{k}=w_{k-1}$,
so the algorithm terminates naturally.
\item \textbf{Extension to Vectors.} For $w\in\mathbb{R}^{d}$ the same proof
holds coordinate‑wise, yielding a diagonal matrix $G_{k}=\operatorname{diag}
\bigl(\sqrt{\sum_{i=1}^{k}g_{i,1}^{2}},\dots,\sqrt{\sum_{i=1}^{k}g_{i,d}^{2}}\bigr)$.
The regret bound becomes $\sum_{j=1}^{d}D_{j}\sqrt{\sum_{k}g_{k,j}^{2}}$.
\end{itemize}

\end{document}